\chapter{Exact Arithmetic}\label{ch:q}

The first thing a computer-algebra system must get right is arithmetic, and the first thing to
relearn is why that is not trivial. A calculator says $0.1 + 0.2 = 0.30000000000000004$; a
mathematician says $\frac{1}{10} + \frac{2}{10} = \frac{3}{10}$. The difference is the difference
between a floating-point number and a rational number, and the engine is built on the second.

\section{The rational numbers}

\begin{definition}[Rationals as pairs]
  A rational number is an equivalence class of pairs $(a, b)$ of integers with $b \neq 0$, under
  the relation $(a, b) \sim (c, d)$ iff $ad = bc$. Write the class as $a/b$.
\end{definition}

That $\sim$ is an equivalence relation is the first exercise of the subject and worth doing once
more. Reflexivity and symmetry are immediate. For transitivity, suppose $ad = bc$ and $cf = de$;
then $adf = bcf = bde$, and since $d \neq 0$, $af = be$. The cancellation of $d$ is the step that
uses that the integers are an integral domain: $xy = 0$ implies $x = 0$ or $y = 0$.

Addition and multiplication are defined on representatives,
\[
  \frac{a}{b} + \frac{c}{d} = \frac{ad + bc}{bd}, \qquad
  \frac{a}{b} \cdot \frac{c}{d} = \frac{ac}{bd},
\]
and one checks that the result does not depend on the representative chosen. With these
operations $\mathbb{Q}$ is a \emph{field}: a commutative ring in which every nonzero element has
a multiplicative inverse, namely $(a/b)^{-1} = b/a$.

\begin{theorem}[Reduced form]
  Every rational has a unique representative $a/b$ with $b > 0$ and $\gcd(a, b) = 1$.
\end{theorem}
\begin{proof}
  Existence: given any representative $a/b$, divide both by $g = \gcd(a, b)$ and, if the
  denominator is negative, negate both. Uniqueness: if $a/b = c/d$ are both reduced with positive
  denominators, then $ad = bc$; $b$ divides $ad$ and is coprime to $a$, so $b \mid d$; symmetrically
  $d \mid b$; so $b = d$ and then $a = c$.
\end{proof}

The step ``$b \mid ad$ and $\gcd(a,b) = 1$ implies $b \mid d$'' is Euclid's lemma, and its proof
is Bézout's identity: there are integers $u, v$ with $ua + vb = 1$, so $d = uad + vbd$, and $b$
divides both terms. Bézout's identity is what the Euclidean algorithm computes. So the reduced
form of a fraction is not a convention; it is an algorithm, and it is the algorithm Lean's
\lc{Rat} runs on every operation.

\section{Why floating point is not a field}

A double-precision float is a dyadic rational $m \cdot 2^e$ with $|m| < 2^{53}$, together with a
rounding rule for results that are not representable. The rounding is what breaks the axioms.
Addition is not associative: with $a = 10^{16}$, $b = -10^{16}$, $c = 1$, the sum $(a + b) + c$ is
$1$ while $a + (b + c)$ is $0$, because $b + c$ rounds back to $b$. Distributivity fails
similarly. None of the identities a simplifier uses --- collecting $x + x$ into $2x$, cancelling
$x/x$ --- is exactly true of floats.

The engine therefore does all arithmetic in $\mathbb{Q}$, and keeps a single bit that records
whether a number \emph{arrived} as a decimal. A literal \texttt{2.5} is stored as $5/2$ exactly,
flagged approximate, and printed back as a decimal; the flag is contagious through arithmetic
and means nothing to any theorem.

\begin{minted}{lean}
structure Q where
  val : Rat
  approx : Bool := false

def add (a b : Q) : Q := ⟨a.val + b.val, a.approx || b.approx⟩
def mul (a b : Q) : Q := ⟨a.val * b.val, a.approx || b.approx⟩
def parse (s : String) : Option Q   -- "3", "2.5" (as 5/2, approximate), "1/3"
\end{minted}

\section{Decimal expansions}

A rational printed as a decimal either terminates or repeats, and the relearning is in the
\emph{why}. Dividing $a$ by $b$ by long division, each step produces a remainder in
$\{0, \ldots, b-1\}$; after at most $b$ steps a remainder recurs, and from that point the digits
repeat. Conversely a repeating decimal $0.\overline{d_1 \ldots d_k}$ is $\frac{d_1 \ldots d_k}{10^k - 1}$,
a rational. So the rationals are exactly the numbers with eventually periodic expansions, and
$\sqrt 2$, whose expansion is not periodic (Chapter~\ref{ch:rad}), is not one of them.

When the notebook is asked for a decimal it shows fifteen significant digits, which is what a
double carries, and rounds the exact rational to them. The rounding is a function on
$\mathbb{Q}$, \lc{Q.roundSig}, and its result is again an exact rational: nothing in the engine
ever holds a float, even when it prints one.

\section{The engine's terms}

Before the first rule, the objects it acts on. A term is a tree:

\begin{minted}{lean}
inductive Expr where
  | num    : Q → Expr
  | var    : String → Expr
  | add    : List Expr → Expr
  | mul    : List Expr → Expr
  | pow    : Expr → Expr → Expr
  | fn     : String → List Expr → Expr
  | matrix : List (List Expr) → Expr
\end{minted}

There is no subtraction, division or negation: $a - b$ is $a + (-1)\cdot b$, $a/b$ is
$a \cdot b^{-1}$. Fewer constructors means fewer cases in every proof, and the printer restores the
human notation. Sums and products take lists, because $a + b + c$ is one node, not two, and
because collecting like terms is a statement about a list.

A rule is a partial function on terms together with an explanation and, in the simplifier, a
proof that it makes a measure decrease (Chapter~\ref{ch:wf}):

\begin{minted}{lean}
structure Rule (W : Weights) where
  name : String
  apply : Expr → Option RuleResult
  decreasing : ∀ e r, apply e = some r → measure W r.result < measure W e
\end{minted}

And a rule \emph{means} something through an evaluator. The engine has three, each a restriction of
the next. The integer evaluator is decidable and partial: it gives a value to integer numerals,
variables, sums, products and natural powers, and \lc{none} to everything else.

\begin{minted}{lean}
def eval? (ρ : Env) : Expr → Option Int
  | .num q => q.asInt?
  | .var x => some (ρ x)
  | .add es => evalSum? ρ es
  | .mul es => evalProd? ρ es
  | .pow b e => do let v ← eval? ρ b; let n ← eval? ρ e; if 0 ≤ n then some (v ^ n.toNat) else none
  | .fn _ _ => none
  | .matrix _ => none

def Refines (e e' : Expr) : Prop := ∀ ρ v, eval? ρ e = some v → eval? ρ e' = some v
\end{minted}

The real evaluator, in \file{proofs/}, is total and is where the mathematics of this book mostly
lives:

\begin{minted}{lean}
def evalR (ρ : EnvR) : Expr → ℝ
  | .num q => (q.val : ℝ)
  | .var x => ρ x
  | .add es => sumR ρ es
  | .mul es => prodR ρ es
  | .pow b e => (evalR ρ b) ^ (evalR ρ e)
  | .fn f [a] => applyFn f (evalR ρ a)
  | .fn _ _ => 0
  | .matrix _ => 0

def SemEqR (a b : Expr) : Prop := ∀ ρ, evalR ρ a = evalR ρ b
\end{minted}

A rule is \status{verified} over the reals when \lc{evalR} of its output equals \lc{evalR} of its
input in every environment. The bridge theorem \lc{evalR_of_eval?} says the two evaluators agree
wherever the first is defined, so a theorem over $\mathbb{R}$ is never weaker than one over
$\mathbb{Z}$.

\section{The rule: folding constants}

The first rule collects the numerals of a sum or product into one:

\begin{minted}{lean}
def foldApply : Expr → Option RuleResult
  | .add es =>
    let nums := es.filter isNum
    if 2 ≤ nums.length then
      some ⟨.add (.num (sumQ nums) :: es.filter (fun e => !isNum e)), "Arithmetic on constants …", none, none⟩
    else none
  | .mul es =>
    let nums := es.filter isNum
    if 2 ≤ nums.length then
      some ⟨.mul (.num (prodQ nums) :: es.filter (fun e => !isNum e)), "Arithmetic on constants …", none, none⟩
    else none
  | _ => none
\end{minted}

\lc{sumQ} folds \lc{Q.add} over the numerals. The rule fires only when there are at least two,
which is what makes it a decrease: two nodes become one.

\section{The theorem}

\begin{thmbox}[\lcn{foldConstants_soundR}]
  For every environment $\rho$ and term $e$ on which the rule fires with result $e'$,
  $\lc{evalR}\,\rho\,e = \lc{evalR}\,\rho\,e'$.
\end{thmbox}

The proof has one mathematical content and one bookkeeping content. The mathematics is that the
cast $\mathbb{Q} \to \mathbb{R}$ is a ring homomorphism: $\overline{p + q} = \overline p +
\overline q$ and $\overline{pq} = \overline p \, \overline q$. In Mathlib these are
\lc{Rat.cast_add} and \lc{Rat.cast_mul}, and they are the reason a fold in $\mathbb{Q}$ is a
fold in $\mathbb{R}$. The bookkeeping is that filtering a list into its numerals and its
non-numerals and adding the two parts back gives the original sum: \lc{sumR} respects
\lc{List.filter} and permutation (\lc{sumR_perm}, an induction on \lc{List.Perm}).

\begin{minted}{lean}
theorem sumQ_evalR (ρ : EnvR) (l : List Expr) (hall : ∀ e ∈ l, Expr.isNum e = true) :
    ((sumQ l).val : ℝ) = sumR ρ l
theorem foldConstants_soundR : RuleSoundR foldConstants
\end{minted}

The status is \status{verified}, and the notebook's note for it reads ``exact rational arithmetic;
$\mathbb{Q}$ embeds in $\mathbb{R}$'' --- which is this section in six words.

\begin{trybox}
\begin{minted}{text}
1/3 + 1/6
0.1 + 0.2
2.5 * 4
N(1/3)
\end{minted}
  The first is $1/2$. The second is $0.3$, because the decimals were exact all along and only
  the printing is approximate. The fourth asks for the decimal explicitly.
\end{trybox}

\section{Exercises}

\begin{exercisebox}[Well-definedness]
  Show that $\frac{a}{b} \cdot \frac{c}{d} = \frac{ac}{bd}$ does not depend on the representatives:
  if $(a,b) \sim (a',b')$ and $(c,d) \sim (c',d')$ then $(ac, bd) \sim (a'c', b'd')$.
\end{exercisebox}

\begin{exercisebox}[Bézout by hand]
  Run the extended Euclidean algorithm on $(240, 46)$ and write $\gcd(240, 46)$ as
  $240u + 46v$.
\end{exercisebox}

\begin{exercisebox}[A float counterexample]
  Find three doubles for which $(a \cdot b) \cdot c \neq a \cdot (b \cdot c)$, and explain which
  rounding is responsible.
\end{exercisebox}

\begin{exercisebox}[The bridge]
  \lc{evalR_of_eval?} is proved by induction on the term. Which case needs
  \lc{Real.rpow_natCast}, and why does the integer evaluator refuse negative exponents?
\end{exercisebox}
